\noindent\hspace*{\parindent}The following explains items (3) and (8). The oldest of the present members is numbered 135: the first 134 have all died or left in some other way. For each of these 134 former members, the following has been calculated: the amount of the contributions he paid and, where survivors have received pensions, the amount of the widows' pensions and orphans' pensions. All these payments, however, have been entered at the rates currently in force rather than at their actual amounts, namely 15 thalers for the annual contribution and 200 thalers as the present amount of the fixed pension; in addition, all money paid or received has been discounted to the time when the member concerned joined. For the adjustments that still had to be made to these calculations for those of the 134 members whose widows are still living, I refer to my memorandum of 1845.

Understood in this sense, the totals for all these 134 members are
\begin{center}
\begin{tabular}{l r l}
the sum of all contributions . . . . . .&25506&thalers\\
the sum of all widows' pensions&66771&,,\\
the sum of all orphans' pensions&9604&,,
\end{tabular}
\end{center}
\noindent which gives the following average values: for contributions, 190,35 thalers; for widows' pensions, 498,29 thalers; for orphans' pensions, 71,67 thalers. In 1845, when the uninterrupted sequence of members who had left by then extended only to No. 108, I had carried out a very similar calculation for them. When put into the same form, it had given the following average values: 193,18 thalers; 503,08 thalers; 81,07 thalers. It is satisfying to see from this that the experience accumulated so far is already very suitable for establishing average values and thus provides a valuable basis for estimating future needs. According to the new result, the average ratio of the orphan's pension to the widow's pension is very nearly \(\tfrac{1}{7}\); this was used to establish item (2) (compare the end of p. [165]).

Furthermore, according to the average values, the gradual contributions of each new member, on the one hand, and the pension payments to his survivors, on the other, are equivalent to the sums of 190,35 and 569,96 thalers paid into and out of the fund, respectively, in a single payment at the time of his entry. Thus, if one new member joined each year, it would only be necessary to imagine these two payments repeated each year and then capitalize them, that is, assign to the fund once and for all the capital sum of 5439 thalers as income and charge the other 16285 thalers as expenditure, to take account of the obligations and entitlements of all members joining in future. These sums must, however, still be multiplied by the figure expressing the \SourceLeaf{172}average number of new members joining each year. In the memorandum of 1845 I assumed \(2\tfrac{1}{2}\) for this, without concealing that this element is very uncertain. Having weighed everything, I believed I had to retain the same figure this time too, and this gave the figures entered under items (8) and (3) in the balance account.
